\section{Introduction}
\label{sec:intro}

Capacity utilization is among the most widely reported macroeconomic
ratios and among the least identified. Central banks read it as a gauge
of inflationary pressure. Growth models treat it as the variable that
absorbs the gap between accumulation and effective demand. A four-decade
heterodox debate over whether firms converge on a normal rate has been
conducted almost entirely on the assumption that the rate is observed
\citep{Kurz1986, BleckerSetterfield2019}. It is not. The numerator of
the ratio---realized output---is measured. The denominator---productive
capacity---is not. Every claim about utilization is therefore a claim
about an object no statistical agency records, resting on whatever
convention supplied the denominator.

The utilization ratio has been supplied in three ways, and all three
fail identically. The first takes utilization as a \emph{given} macro
ratio---an observed series to be used as a regressor or a policy
input---when its denominator was never observed. The second anchors it
to an \emph{inherited benchmark}: an engineering maximum or a
survey-reported normal rate imported from outside the accounting
framework that is supposed to give the ratio meaning. Managerial
surveys do not support the premise this requires
\citep{Singh2022, Nikiforos2016}. The third recovers utilization as an
\emph{empirical residual}: whatever remains after output has been
fitted to capital by a filter or a production function. This third
treatment is where most current empirical work sits, and the objection
to it is structural. Fitting a trend through output conditional on
capital sets the capacity path to move one-for-one with
accumulation---$\theta = 1$ by construction---and then defines
utilization as the deviation from a benchmark that assumed away the
very relation needed to locate it. The three treatments are not three
difficulties but one: each supplies the denominator by convention, and
each conceals the transformation relation that determines it.

Chapter~1 demonstrated that the bivariate output--capital relation
fails system-level cointegration and that a stable long-run capacity
path requires the rate of exploitation to enter the cointegrating
vector, yielding a transformation elasticity below unity
($\hat{\theta} < 1$) indicative of structural overaccumulation. That
finding established \emph{that} the output--capital relation is
distributionally conditioned. It did not identify \emph{how}
distribution conditions the conversion of accumulated capital into
productive capacity, nor did it extend the analysis beyond the
closed-economy core. This chapter supplies both.

I estimate the transformation elasticity $\theta \equiv \partial \ln
Y^p / \partial \ln K$ as a state-dependent cointegrating relation.
For the United States over 1931--2024, I use the Integrated Modified
OLS estimator of \citet{VogelsangWagner2024} with fixed-$b$ inference
and endogenously dated structural breaks \citep{BaiPerron2003};
distribution enters through effective wage pressure. For Chile over
1943--2010, the balance-of-payments constraint enters as a threshold
that switches the choice of technique. The estimates locate $\theta$
well below unity in both economies. In the United States the elasticity
moves from $0.37$ in 1960 to $0.43$ in 2024 as effective wage
pressure falls with the diversion of investment into intangibles;
utilization declines from its 1944 wartime peak to $93.8\%$ in 2024.
In Chile a composite balance-of-payments threshold at
$\hat{\gamma} = 0.425$ switches the mechanization response off: the
wage-share elasticity of the capital composition gap falls from $+3.29$
to $+0.14$ once foreign exchange binds \citep{Hansen2000}. The
center--periphery comparison establishes that the same transformation
relation is mediated by distributive conflict in the center and
truncated by the balance-of-payments constraint in the periphery.

Two literatures have circled this problem without closing it. The
induced-innovation tradition following \citet{Kennedy1964} made
distribution govern the \emph{direction} of technical change along a
frontier whose \emph{scale} stayed exogenous; its own later assessments
concede that it explains the composition of innovation while remaining
close to silent on its intensity \citep{TavaniZamparelli2018}. The
measurement literature has established convincingly that reported
capacity is a construct rather than an observation
\citep{Shaikh2016, Nikiforos2016}, but has stopped short of supplying
the relation that would replace it. \citet{Okishio1961, Okishio2022}
supplies half of that relation---decentralized investment produces
disproportionality as a structural norm, not a friction---and
\citet{Vidal2022, Vidal2019} supplies the other half, showing that the
firm converting capital into usable capacity is internally divided
over control and cooperation and settles into configurations that
reproduce production without resolving the conflict. Neither framework
was carried to the point where the transformation elasticity could be
estimated. \citet{Basu2022, Basu2026} closes part of the gap by tying
the profitability outcome of viable technical change to the
labor-market closure, which maps onto historical accumulation regimes;
and for the periphery, \citet{Prebisch1981} and \citet{Kaldor1959}
established the external ceiling without connecting it to the choice of
technique that runs into it.

This chapter makes three contributions against that background. First,
it reconstructs capacity utilization as a \emph{derived} quantity,
identified downstream of a transformation elasticity that is itself
estimated rather than assumed, extending the distributional conditioning
established in Chapter~1 from a cointegration-survival result to a
structural elasticity with a historical trajectory. Second, it supplies
an econometric identification of that elasticity that does not
presuppose the relation it recovers, using a cointegrating
specification in which distribution enters as a state-dependent
interaction and structural change is dated endogenously. Third, it
extends the relation to the periphery by letting the balance-of-payments
constraint enter as a \emph{threshold} that switches the choice of
technique, rather than as a control variable added to it.

The argument runs as follows. The transformation relation cannot be
posited without a theory of the unit that performs the transformation,
which is what \Cref{sec:2:divided_firm} supplies by replacing the
unified firm with decentralized accumulation and internal division.
\Cref{sec:3:analytical_framework} then asks what has to be true of the
choice of technique for that transformation to have a determinate
elasticity, and derives $\theta$ under homogeneous capital, under
heterogeneous capital with a mechanization gap, and under an external
constraint that can bind. \Cref{sec:4:empirics} confronts the
resulting specifications with reconstructed capital stocks for both
economies and reports what the estimates say about the historical path
of $\theta$ and $\mu$. \Cref{sec:5:comparative_historical_analysis} sets the two histories
against each other, and \Cref{sec:conclusion} states what the
reconstruction settles and what it leaves open.